\ch{Chapter 1: Introduction}\label{chap:Chapter_1}



When threat actors began exploiting zero-day vulnerabilities in Ivanti Connect Secure and Ivanti Policy Secure appliances in December 2023, security teams across multiple agencies worked under acute time pressure. Defenders sought to identify every affected component and to understand what attackers were already exploiting~\parencite{Ivanti_Threat_Impact}. Yet they had no comprehensive understanding of how vulnerabilities in one system could cascade across interconnected infrastructure and national security systems. When the Cybersecurity and Infrastructure Security Agency issued Emergency Directive 24-01 in January 2024, agencies lacked ready visibility into the affected configurations and their dependencies. While they worked to establish that visibility, adversaries retained the initiative~\parencite{CISA_Emergency_2023}.
% Verified: CISA issued ED 24-01 (Mitigate Ivanti Connect Secure and Ivanti Policy Secure Vulnerabilities) on January 19, 2024, covering CVE-2023-46805 and CVE-2024-21887; Volexity detected exploitation beginning the second week of December 2023 and published its report January 10, 2024. Per cisa.gov ED 24-01 pages and volexity.com, July 2, 2026. The CISA_Emergency_2023 bib entry content (title, ED number, January 2024) is correct; only its citation key is misleadingly named.
This operational failure is not an isolated incident. It exemplifies a structural condition of modern practice: enterprises now change their information systems faster than their documentation can describe them, leaving defenders to act on records that no longer match reality.

The consequences of such failures extend beyond the immediate organizations affected. When defenders cannot comprehend the cascading impacts of compromise, they communicate risk less clearly to organizational leadership, they lose the ability to prioritize remediation by actual impact severity, and coordinating a defense across organizational boundaries becomes impractical. When defenders cannot understand system interdependencies, incidents that might otherwise be contained become enterprise-wide crises. Security teams then shift to reactive responses rather than proactive defense, spending effort on manual discovery. Model-based approaches are designed to do that same work far more efficiently, a goal the Department of Defense set out in its 2018 Digital Engineering Strategy~\parencite{DoD_DE_Strategy_2018}.

Information Assurance, as codified by the National Institute of Standards and Technology, encompasses those measures that protect and defend information and information systems by ensuring their availability, integrity, authentication, confidentiality, and non-repudiation~\parencite{Force_2020}. Cybersecurity is an operational component within this broader discipline, concentrating on the protection of information systems from unauthorized access, use, disclosure, disruption, modification, or destruction. Throughout this dissertation, the term \textit{Information Assurance} denotes the broader discipline encompassing security policy, risk management, compliance verification, and protective measures. Cybersecurity, by contrast, refers specifically to the technical and operational dimensions of protecting systems from cyber threats. The two terms are not used interchangeably: Information Assurance is the broader governance and assurance framework, while cybersecurity addresses the protective mechanisms and threat responses that operate within it.

Terminological precision bears operational consequences. Treating Information Assurance as a synonym for cybersecurity invites organizations to under-invest in the governance, documentation, and architectural foundations that effective cybersecurity operations rest on. Failing to maintain accurate system documentation, for example, represents an Information Assurance shortfall that manifests as cybersecurity operational degradation. Literature from defense and aerospace contexts suggests that Digital Engineering may address both dimensions: the governance and documentation requirements of Information Assurance and the operational visibility requirements of cybersecurity defense. Whether IT and Information Assurance professionals recognize this potential is a central question of this research.

This chapter examines the challenges organizations face in implementing Information Assurance practices and managing Information Technology (IT) service delivery. It introduces Digital Engineering as a disciplinary approach to gaps that persist despite mature frameworks, among them the National Institute of Standards and Technology (NIST) Risk Management Framework (RMF)~\parencite{Force_2020}, the Information Technology Infrastructure Library (ITIL)~\parencite{Cannon_AXELOS_2013}, and the Unified Architecture Framework (UAF)~\parencite{OMG_UAF_2022}.

\section{A Note on Institutional Nomenclature}

A point of nomenclature warrants clarification at the outset, because it recurs throughout this dissertation and concerns the very institution from which the methods studied here originate. The United States defense establishment traces its origins to the Department of War, established in 1789 to administer primarily the Army. The National Security Act of 1947 created the National Military Establishment, unifying the War and Navy Departments together with the newly independent Air Force. That body was renamed the Department of Defense in 1949. The institution has now entered a transition back toward its original title. In September 2025, Executive Order 14347 authorized the use of ``Department of War'' as a secondary title for the institution and its senior officials in official correspondence, public communications, and other non-statutory contexts~\parencite{EO_14347_2025}. Only an Act of Congress can alter the statutory name of an executive department. The legally established name therefore remains the Department of Defense while the secondary designation is in active use and legislation to make the change permanent is pending. This dissertation accordingly uses the historical names when referring to documents and events of their own eras and ``Department of War'' when referring to the institution as a present-day actor.

Concretely, the dissertation refers to the institution as the Department of War (DoW) wherever it appears as a present-day actor, for example in describing its Digital Engineering Strategy, while preserving published titles exactly as issued. The Department of Defense Architecture Framework remains DoDAF, Department of Defense Instruction 5000.97 remains DoDI 5000.97, and bibliographic entries credit the Department of War as institutional author while keeping each document's published title verbatim. The convention keeps the work current with the institution's evolving identity, which matters because the Department of War is the principal originator of the Digital Engineering practices this research examines. It also preserves the scholarly record: cited titles remain exactly as published and therefore discoverable. Where the abbreviation DoD appears in later chapters it does so only as an inseparable element of a proper name, such as DoDAF or DoDI, and never as a reference to the institution itself.

\section{Current State of Information System Management}

Organizations today operate within an environment defined by rapid technological evolution and escalating system complexity. Cloud computing, microservices architectures, Internet of Things (IoT) devices, and operational technology have converged to produce complex, interdependent systems that overwhelm traditional approaches to both Information Assurance and IT service delivery~\parencite{Benbya_Nan_Tanriverdi_Yoo_2020}. These advances deliver undeniable operational benefits, but they carry substantial costs in system visibility, security control implementation, configuration management, and service delivery coordination.

Technological change continues to accelerate. Organizations that required months to deploy new capabilities a decade ago now deploy changes continuously through automated pipelines. This acceleration benefits operational agility but strains the documentation and verification processes on which Information Assurance depends: static documentation designed for quarterly or annual update cycles cannot maintain accuracy when systems change hourly, a mismatch examined later in this chapter.

Enterprise information systems now routinely span multiple technology domains: cloud-based infrastructure and services, on-premises data centers, edge computing environments, operational technology networks, and mobile and remote access systems. This technological heterogeneity generates persistent challenges in maintaining security visibility, implementing consistent protection mechanisms, and delivering reliable IT services across disparate environments. Teams cannot track system dependencies with confidence, configuration management falters, and security control implementation proceeds inconsistently across heterogeneous platforms~\parencite{Bokan_Santos_2021}.

Beyond technical complexity, the challenge extends to organizational design. Structures that evolved to manage discrete technology domains now impede the integrated visibility that modern environments require. Security teams operate separately from IT operations teams, while cloud architects work independently of network engineers. Application developers deploy services without understanding infrastructure dependencies. This organizational fragmentation mirrors and reinforces the technical fragmentation that undermines both Information Assurance and IT service delivery effectiveness.

\subsection{Information Assurance Practice}

Information Assurance practice has evolved in response to the complexities of modern enterprise environments. Central to this evolution, the NIST Risk Management Framework provides a structured, disciplined approach for managing security and privacy risk that organizations can apply across diverse information systems~\parencite{Force_2020}. RMF establishes a lifecycle approach to security through seven iterative steps: prepare, categorize, select, implement, assess, authorize, and monitor. Now foundational for federal agencies, the framework finds increasing adoption among organizations operating national security systems and private enterprises seeking systematic approaches to Information Assurance.

The RMF represents a substantial advancement over earlier compliance-focused approaches that treated security as a point-in-time certification rather than a continuous process. Its emphasis on continuous monitoring and ongoing authorization reflects recognition that security postures change constantly as systems evolve, threats emerge, and organizational requirements shift. Implementing continuous monitoring effectively requires capabilities that most organizations lack: real-time visibility into system configurations, automated assessment of control effectiveness, and dynamic risk calculation based on current rather than documented system states.

Additional Information Assurance lifecycle frameworks exist, including ISO 31000~\parencite{for_2018}, the NIST Cybersecurity Framework~\parencite{Cybersecurity_Framework_2013}, and COBIT~\parencite{ISACA2019COBIT}, among others. These frameworks provide alternative approaches to the RMF\@. However, they share common challenges in maintaining accurate documentation, ensuring visibility into system states, and coordinating security efforts across organizational boundaries. This research focuses on the NIST Risk Management Framework, a choice that reduces complexity by avoiding direct comparison among multiple frameworks while still addressing challenges common to all.

Security control selection is a key RMF activity. Federal information systems and national security systems typically use NIST Special Publication 800-53 Revision 5 as the authoritative catalog of security controls~\parencite{Force_2020}. Organizations operating outside federal requirements may employ alternative frameworks for control selection, including ISO/IEC 27001~\parencite{ISO27001-2023}, the NIST Cybersecurity Framework, or industry-specific standards. This research focuses on NIST 800-53 Revision 5 given its applicability to federal and national security contexts, though the underlying principles extend to organizations employing other control frameworks.

NIST Cybersecurity Framework (CSF) 2.0, released in February 2024, provides a voluntary, risk-based structure that complements the RMF by organizing cybersecurity outcomes into six core Functions: Govern, Identify, Protect, Detect, Respond, and Recover~\parencite{NIST_CSF_2_0_2024}. Version 2.0 added the Govern Function to address cybersecurity governance at the organizational level. This addition recognizes that leadership engagement, risk strategy, and policy oversight are prerequisites for effective cybersecurity rather than ancillary concerns. Each Function decomposes into Categories and Subcategories that map to specific outcomes. CSF 2.0 maintains explicit alignment with SP~800-53 controls through Informative References, enabling organizations to trace CSF outcomes to specific control implementations. CSF 2.0 broadens its scope beyond critical infrastructure to encompass all organizations regardless of sector or size. In this expanded form, it serves as a bridging framework. Organizations subject to RMF compliance can use CSF to communicate risk posture to non-technical stakeholders. Organizations not subject to federal requirements can adopt CSF as a starting framework that maps naturally to SP~800-53 controls should compliance requirements evolve. CSF 2.0's emphasis on governance, supply chain risk management, and continuous improvement aligns with Digital Engineering principles, particularly the authoritative source of truth concept and digital thread traceability. To the author's knowledge, no published research has examined this alignment.

Beyond federal systems, an expanding compliance landscape governs nonfederal organizations that process, store, or transmit Controlled Unclassified Information (CUI) on behalf of government agencies. NIST Special Publication 800-171 Revision 3, finalized in May 2024, establishes security requirements for protecting CUI in nonfederal systems and organizations~\parencite{NIST_SP_800_171_R3_2024}. Revision 3 restructured the framework from fourteen to seventeen security requirement families, adding Planning, System and Services Acquisition, and Supply Chain Risk Management. It also aligned its control structure directly with SP~800-53 Revision 5 as the single authoritative source. % Verified: the final SP 800-171r3 (May 2024) contains 97 security requirements across 17 families, per the official publication (doi:10.6028/NIST.SP.800-171r3), July 2, 2026. The counts "266 control items" and "49 ODPs" appear nowhere in the final publication (Appendix D lists ODPs but states no total, and third-party ODP counts conflict), so both numbers were dropped.
SP~800-171's 97 security requirements incorporate Organization-Defined Parameters that provide implementation flexibility while enabling automated compliance assessment through machine-readable formats such as the Open Security Controls Assessment Language (OSCAL)~\parencite{NIST_OSCAL_2023}. CSF 2.0 maps to SP~800-171 requirements through shared SP~800-53 lineage: organizations implementing SP~800-171 controls simultaneously address CSF Subcategory outcomes in the Identify, Protect, Detect, Respond, and Recover Functions, creating a unified compliance posture across both frameworks.

NIST Special Publication 800-172 supplements the SP~800-171 baseline with thirty-five enhanced security requirements designed to protect CUI associated with critical programs and high-value assets against Advanced Persistent Threats~\parencite{NIST_SP_800_172_2021}. Operating on a three-pillar defense strategy that comprises penetration-resistant architecture, damage-limiting operations, and cyber resiliency and survivability, SP~800-172 demands capabilities that manual processes and static documentation struggle to sustain. Enhanced requirements span access control, configuration management, identification and authentication, incident response, risk assessment, situational awareness, and system protection. A Revision 3 update currently in development expands scope from confidentiality protection alone to encompass confidentiality, integrity, and availability, further increasing documentation and verification demands.

Operationalizing these NIST requirements within the defense industrial base, the Cybersecurity Maturity Model Certification (CMMC) 2.0 framework establishes a three-tiered certification structure~\parencite{DoD_CMMC_2024}. Level 1 requires the fifteen basic safeguarding requirements of Federal Acquisition Regulation (FAR) clause 52.204-21 for Federal Contract Information protection, verified through annual self-assessment.
% Verified: the final CMMC rule (32 CFR Part 170, published October 15, 2024) sets Level 1 at the 15 security requirements of FAR 52.204-21 (48 CFR 52.204-21(b)(1)), assessed by annual self-assessment. Per ecfr.gov 32 CFR 170 and federalregister.gov 2024-22905, July 2, 2026.
Level 2 mandates compliance with the 110 requirements of SP~800-171 Revision 2, assessed through triennial third-party evaluation by CMMC Third-Party Assessment Organizations. Level 3 adds twenty-four enhanced requirements drawn from SP~800-172, assessed by the Defense Industrial Base Cybersecurity Assessment Center. Phase 1 implementation began with the December 2024 effective date of 32 CFR Part 170, with full implementation across all contracts involving Federal Contract Information or Controlled Unclassified Information expected by November 2028. CMMC marks a fundamental shift from self-attestation to verified compliance, imposing documentation and evidence requirements that intensify the challenges organizations already face in maintaining accurate, current security documentation. This expanding compliance landscape extends from CSF 2.0's voluntary governance framework through SP~800-171 and SP~800-172's CUI protection requirements to CMMC's third-party verification mandate. Across its full breadth, organizations confront escalating documentation, traceability, and verification demands that Digital Engineering capabilities are designed to address.

Selecting appropriate security controls depends on accurately understanding the systems being protected, their operational context, their interconnections with other systems, and their role within the broader enterprise architecture. Control selection that proceeds from inaccurate system understanding produces security postures that address documented rather than actual risk. This disconnect between documentation and reality is a structural challenge that persists regardless of which control framework an organization employs.

Implementing the RMF effectively presents documented challenges in complex technological environments. Organizations must categorize information systems based on potential impact, select appropriate security controls from comprehensive catalogs, implement those controls across diverse platforms, assess control effectiveness, obtain authorization decisions, and maintain continuous monitoring throughout the system lifecycle. Each step demands accurate, current information about system configurations, security control implementations, and operational states. Traditional documentation approaches struggle to maintain such information.

The continuous monitoring requirement deserves particular attention because it exposes the limitations of document-centric approaches most directly. Continuous monitoring as envisioned by the RMF requires ongoing awareness of security-relevant system changes, automated assessment of security posture impacts, and timely reporting to authorizing officials. Organizations attempting to implement continuous monitoring through manual processes discover that the labor required exceeds available resources. Those pursuing automation instead discover that they lack the authoritative system models and configuration baselines that automation requires.

\subsection{IT Service Management Practice}

IT service management (ITSM) has matured through frameworks designed to ensure reliable service delivery across the enterprise. The Information Technology Infrastructure Library provides comprehensive guidance for aligning IT services with business needs through structured processes for service strategy, service design, service transition, service operation, and continual service improvement~\parencite{Cannon_AXELOS_2013}. ITIL emphasizes configuration management, change management, and service asset management as core capabilities on which effective IT service delivery depends.

ITIL's evolution from Version 3 to ITIL 4 reflects recognition that service management practices must adapt to cloud computing, DevOps practices, and agile delivery models. ITIL 4 introduces the Service Value System concept, emphasizing flexibility and continuous improvement over rigid process compliance. Yet the dependencies on accurate configuration information and effective change coordination persist regardless of which ITIL version organizations adopt. Even the Service Value System cannot create value if the underlying information about services, configurations, and dependencies remains inaccurate or incomplete.

Configuration management within the ITIL framework requires organizations to maintain accurate configuration management databases documenting configuration items, their attributes, and their relationships. Change management processes depend on accurate configuration information to assess change impacts and coordinate modifications across interconnected systems. Service asset management extends these capabilities to encompass the full lifecycle of IT assets from acquisition through retirement. These interconnected processes provide structure for managing complex IT environments, but they depend on the accuracy and currency of underlying information, an accuracy that organizations consistently fail to achieve.

The relationship between configuration management and change management illustrates how documentation failures compound. Change management processes assess proposed changes against documented configurations and relationships. When documentation is incomplete, change assessments miss dependencies that exist in operational systems. Changes approved on the basis of incomplete assessments cause unintended impacts. Those impacts require emergency changes to address. Emergency changes bypass change management processes, further degrading documentation accuracy. This cycle perpetuates itself, progressively undermining both configuration management and change management effectiveness.

\subsection{Challenges in Current Practice}

Traditional documentation approaches and manual tracking methods prove increasingly inadequate for capturing and managing the complexity of modern information systems. Paper-based security documentation, static network diagrams, and periodic compliance assessments fail to reflect the dynamic nature of contemporary enterprise environments. IT service management practices that rely on manual configuration tracking and change coordination struggle to maintain accuracy and timeliness in environments characterized by continuous deployment and rapid change cycles.

A structural challenge emerges: the frameworks are sound and the practitioners dedicated, yet the limitation lies in the mismatch between the documentation velocity that manual processes can sustain and the operational velocity that modern enterprise environments demand. The evidence reviewed below suggests that process improvement and additional staffing alone cannot close this gap. Closing the gap requires a paradigm shift from document-centric to model-centric practices, precisely the shift that Digital Engineering is designed to provide.

Research documents pervasive failures across both Information Assurance and IT service management domains. Industry analyst estimates indicate that 80\% of Configuration Management Database (CMDB) implementations fail to deliver intended value~\parencite{Gartner_CMDB_2019}. Studies find that organizations can monitor only 66\% of their IT environments, leaving 34\% unmonitored~\parencite{IDC_Exabeam_2023}. Shadow IT, defined as technology acquired or deployed outside official governance, now represents 30--40\% of enterprise IT spending, creating assets invisible to documentation efforts~\parencite{Gartner_Shadow_IT_2022}. The mean time to identify security breaches averages 204 days, reflecting the visibility gaps that impair threat detection~\parencite{IBM_Ponemon_2024}.

These statistics reveal systemic limitations of document-centric approaches rather than isolated organizational shortcomings. Organizations that invest heavily in documentation still experience these failures. Organizations with mature governance processes still discover undocumented systems and unknown dependencies. Such consistency of failure across diverse organizations identifies the document-centric approach itself, however well executed, as the source of the problem.

A common pattern unites these failures: organizations cannot maintain accurate, current documentation of their information systems using traditional approaches. Modern IT environments change at rates exceeding the capacity of manual documentation processes. Static artifacts become obsolete before completion. Configuration databases diverge from operational reality. Security documentation describes intended states rather than actual implementations. Such a documentation-reality gap undermines every process that depends on accurate system information, a category that includes nearly all Information Assurance and IT service management activities.

Chapter 2 examines these visibility and documentation failures in depth, synthesizing peer-reviewed research and industry analysis to establish the evidence base for why traditional practices fail. Documentation challenges reflect multiple converging factors: organizational silos fragmenting visibility, complexity exceeding human documentation capacity, manual processes unable to match rates of change, and technical debt accumulating in documentation domains.

\section{Digital Engineering as Potential Solution}

The preceding section established that these failures persist despite organizational effort, which reframes the problem as one of approach: the method itself falls short even where execution is diligent. Digital Engineering answers that reframed problem by changing the unit of record from the document to the model: it is a systematic approach to designing, developing, and managing complex systems through integrated digital models and data-driven processes~\parencite{DigitalEngineeringInformationExchange}. The discipline originated in the defense and aerospace sectors, where organizations including the United States Department of War (DoW)~\parencite{DoD_Engineering_2022}, the National Aeronautics and Space Administration (NASA)~\parencite{Engineer}, and the International Council on Systems Engineering (INCOSE)~\parencite{DigitalEngineeringInformationExchange} have formalized it through authoritative guidance. While these practices emerged primarily from physical systems engineering, the underlying principles offer capabilities of documented value for Information Technology and Information Assurance domains.

The defense and aerospace origins of these practices carry significance beyond historical interest. Both sectors developed them to address challenges structurally similar to those confronting enterprise IT and Information Assurance: complex interdependent systems, stringent compliance requirements, mission-critical operations, and the need to maintain comprehensive visibility across extended lifecycles. Solutions that proved effective for managing combat aircraft development and spacecraft missions may serve equally well for managing enterprise information systems and security postures. That transfer remains untested, and this dissertation measures professional perception of its potential rather than the transfer itself. Measured evidence supports this potential: \citeauthor{Rogers_Mitchell_2021} documented an 18\% improvement in systems engineering efficiency and a 9\% reduction in defects following MBSE adoption on a complex system-of-systems program~\parencite{Rogers_Mitchell_2021}. Yet such measured evidence remains rare, and no comparable data exist for enterprise IT or Information Assurance contexts.

Four pillars support Digital Engineering: Model-Based Systems Engineering (MBSE), digital threads (the authoritative traceability that connects system lifecycle artifacts), digital twin technology, and Product Lifecycle Management (PLM). Understanding these pillars provides context for examining how Digital Engineering practices might address the challenges identified in current Information Assurance and IT service management practice.

Integrating these pillars distinguishes Digital Engineering from isolated tool adoption. Organizations might implement modeling tools without achieving Model-Based Systems Engineering, or deploy digital twin capabilities without establishing digital thread traceability. Digital Engineering's value emerges when these capabilities function together as an integrated approach, an integration that current Information Assurance and IT service management practices generally do not provide.

\subsection{Model-Based Systems Engineering}

Model-Based Systems Engineering marks a paradigm shift from document-centric practices to model-centric approaches for system development and management~\parencite{Hutchison_2020}. Rather than relying primarily on textual descriptions and static diagrams, MBSE employs formal, executable models that capture system architecture, behavior, requirements, and relationships in structured, machine-readable formats. These models serve as authoritative sources of truth that can be analyzed, simulated, and validated throughout the system lifecycle~\parencite{Hutchinson_2021}.

Distinguishing document-centric from model-centric approaches matters because the contrast defines the fundamental transformation that Digital Engineering proposes. Document-centric approaches produce artifacts such as diagrams, specifications, and procedures that describe systems. These artifacts require human interpretation, cannot be automatically validated for consistency, and provide no mechanisms for maintaining currency as systems evolve. When an organization produces a network diagram, that diagram captures system state at the moment of creation. Every subsequent change renders the diagram incrementally less accurate. Likewise, a security plan describes intended security posture only at the time of writing, and every subsequent system modification creates potential divergence between the plan and operational reality. Model-centric approaches produce executable representations that can be queried, analyzed, and validated automatically. As models change, dependent artifacts can update automatically, and when engineers evaluate proposed changes, models enable impact analysis that documents cannot provide. Maintaining consistency automatically addresses the velocity mismatch between documentation processes and operational change rates described earlier in this chapter.

Formal modeling languages underpin MBSE implementations and provide the technical foundation for these capabilities. SysML, the Systems Modeling Language standardized through the Object Management Group, extends the Unified Modeling Language (UML) with constructs specifically designed for systems engineering~\parencite{Friedenthal_Moore_Steiner_2014}. Where UML focuses on software system modeling, SysML adds requirements modeling, parametric analysis, and physical system representation that UML lacks. SysML defines nine diagram types organized into four diagram categories (requirements, structure, behavior, and parametrics) that span stakeholder requirements, system architecture and composition, operational dynamics, and constraint-based quantitative analysis (detailed in Chapter~2). This taxonomy provides a comprehensive vocabulary for representing complex systems from multiple perspectives simultaneously, a capability that directly supports the multi-stakeholder visibility requirements of both Information Assurance governance and IT service management coordination.

SysML's evolution from 1.x to the emerging v2 specification addresses limitations that impeded broader adoption. Version 2 introduces a redesigned language architecture emphasizing improved precision, expressiveness, and usability. A standardized textual notation accompanies the graphical notation, enabling model creation through text-based interfaces more familiar to software developers and IT professionals. SysML v2 also introduces a standardized Application Programming Interface (API) for tool interoperability, addressing a persistent criticism that MBSE tool ecosystems require vendor-specific integrations. Improved interoperability reduces the vendor lock-in concerns that organizations evaluating MBSE adoption frequently cite. For enterprise IT contexts, the textual notation and API standardization lower adoption barriers by aligning MBSE practices more closely with the text-based configuration and automation workflows that IT professionals employ daily.

The MBSE tool ecosystem encompasses both commercial and open source platforms. Commercial tools, including Dassault Systemes Cameo Systems Modeler and IBM Engineering Systems Design Rhapsody, provide comprehensive modeling environments with enterprise support, training resources, and integration capabilities developed through decades of aerospace and defense application. Open source alternatives have emerged primarily through the Eclipse Foundation's modeling ecosystem. These include Eclipse Papyrus for UML and SysML modeling~\parencite{Eclipse_Papyrus_2024}; Capella's Arcadia-based environment, which Thales developed and organizations have deployed across a variety of industrial contexts~\parencite{Capella_MBSE_2024, Roques_2016}; and the web-based SysON implementation of the emerging SysML v2 specification~\parencite{SysON_2025}.
% Verified: Roques (ERTS 2016, hal-01258014) documents Capella as "successfully deployed in a wide variety of industrial contexts," developed through Arcadia at Thales, with pilots at Airbus, Airbus DS, and Areva; confirmed via HAL API, July 2, 2026. New bib entry Roques_2016 reported for central merging.
These open source tools reduce economic barriers to MBSE adoption. However, their documentation, community resources, and application examples focus predominantly on traditional systems engineering domains. Organizations seeking to apply MBSE tools for enterprise IT or Information Assurance purposes must adapt without domain-specific guidance, an adaptation that remains unexplored in the academic literature.

Architecture frameworks provide the structural foundation for organizing MBSE implementations. Within MBSE practice, architecture frameworks define the viewpoints, views, and model elements that architects employ to represent systems. UAF, discussed in detail in a subsequent subsection, provides the most comprehensive architecture framework for MBSE implementations spanning both defense and commercial domains. Its multi-viewpoint approach aligns naturally with the needs of organizations managing information systems that must satisfy both Information Assurance requirements and IT service delivery objectives. Architects can represent security controls, operational processes, service dependencies, and resource allocations within a single integrated model environment.

For information systems, MBSE principles enable organizations to create formal models of their IT infrastructure, security architectures, and service delivery processes. These models capture the static configuration of systems and, beyond it, the dynamic relationships between components, the flow of information through the enterprise, and the dependencies that affect both security postures and service delivery. Model-based approaches provide enhanced visibility into system complexity, enable automated analysis of security implications for proposed changes, and support more effective planning for IT service delivery requirements. Integrating MBSE with established frameworks such as the NIST RMF and ITIL enables organizations to maintain living models that reflect both security control implementations and service configuration states. However, achieving this integration requires modeling language extensions and domain-specific profiles that researchers have yet to develop or validate for enterprise IT contexts. SysML provides excellent support for modeling physical systems with requirements, behaviors, and structures, but modeling IT services, network configurations, and security controls requires adaptation or extension that practitioners must develop independently. The absence of standardized approaches for modeling enterprise IT in SysML is a technical barrier to adoption that accompanies the awareness and perception barriers this research investigates.

\subsection{Digital Threads}

Digital threads constitute authoritative traceability: the verified, bidirectional connections between requirements, design elements, implementation artifacts, and validation activities that persist throughout a system's lifecycle~\parencite{Singh_Willcox_2018}. The term describes the connections that link authoritative sources---including Model-Based Systems Engineering models, requirements management systems, configuration management databases, and Product Lifecycle Management repositories---into a unified, navigable body of system information. Digital threads ensure that organizations can track how requirements flow through the development and implementation process, identify which system components implement specific capabilities, and verify that implemented solutions satisfy intended requirements. Unlike traditional documentation approaches, digital threads are designed to maintain verified relationships that remain current as systems evolve~\parencite{Zhang_DigitalThread_2024}.

Authoritative traceability merits close attention because it addresses a fundamental limitation of current Information Assurance and IT service management practice. Traditional traceability attempts to maintain connections through manual cross-references, requirements matrices, and documentation linkages. These manual traceability mechanisms require constant maintenance, degrade as systems evolve, and provide no automated verification of consistency. A requirements traceability matrix that maps security controls to implementation artifacts represents a snapshot of intended relationships at the time of creation. Every subsequent change to either controls or implementations creates potential inconsistency that manual processes may not detect. Digital threads establish traceability through model relationships that can update automatically as models change. Queries against digital thread repositories return current rather than historical information. Impact analyses traverse digital thread connections to identify affected components throughout the system architecture. Maintaining currency automatically distinguishes digital threads from the manual traceability approaches that current compliance practices employ.

The Department of Defense's 2018 Digital Engineering Strategy positions the digital thread as essential infrastructure for achieving the authoritative source of truth: the single, trusted baseline of system information from which all stakeholders operate~\parencite{DoD_DE_Strategy_2018}. As a concept, the authoritative source of truth addresses a persistent challenge in both defense acquisition and enterprise IT management: when multiple documentation artifacts describe the same system, discrepancies inevitably emerge, and no definitive mechanism exists for determining which artifact accurately reflects operational reality. Digital threads resolve this challenge by establishing verified connections between all system artifacts, ensuring that changes propagate consistently and that queries return authoritative rather than potentially obsolete information. The Systems Engineering Research Center (SERC) has produced research supporting digital thread implementation, including the Enterprise System-of-Systems Model for Digital Thread Enabled Acquisition, which establishes architectural patterns for maintaining traceability across complex, multi-system environments~\parencite{SERC_Digital_Thread_2018}. DoDI 5000.97 codifies these concepts into mandatory requirements, directing defense programs to maintain digital thread capabilities throughout the acquisition lifecycle~\parencite{DoDI_5000_97_2023}.

For Information Assurance practice, digital threads address gaps in current RMF implementation. The RMF requires organizations to select security controls, implement those controls, and assess their effectiveness throughout the system lifecycle. Digital threads enable organizations to trace security requirements from categorization decisions through control selection, implementation, and assessment activities, so that policy documents connect to technical configurations, and those configurations connect in turn to assessment evidence, all within a single authoritative chain. This traceability supports the continuous monitoring phase of the RMF by maintaining verifiable connections between security requirements, implemented controls, and compliance artifacts. When a security control implementation changes, digital thread connections enable automated identification of affected requirements, dependent controls, and compliance documentation requiring update. When auditors require evidence of control implementation, digital threads provide navigable paths from requirements through implementation to assessment results without manual evidence compilation.

Within IT service management contexts, digital threads align with ITIL configuration management and change management practices. Organizations can trace service delivery requirements to underlying infrastructure components and configuration items, connecting the CMDB to as-built system documentation and operational baselines. This capability supports more accurate impact assessment for changes and more effective root cause analysis for service disruptions. When a proposed change affects a configuration item, digital thread connections reveal which services depend on that item, which security controls protect it, and which compliance requirements govern it. Maintaining current, verified traceability relationships through digital threads reduces the time and effort required for compliance audits while improving the accuracy of both security assessments and service impact analyses.

Implementing digital threads in practice requires integration infrastructure that connects diverse tools and repositories. In defense contexts, this integration has developed over years of investment and standardization. Enterprise IT environments employ different tool ecosystems---among them IT service management platforms like ServiceNow and BMC, security information and event management systems, cloud management consoles, and configuration automation tools---none of which were designed with digital thread connectivity in mind. Establishing digital thread capabilities in enterprise IT contexts therefore requires integration approaches that bridge systems engineering tools and IT management platforms. This integration challenge is a practical barrier that accompanies the conceptual transfer of digital thread principles from defense to enterprise IT contexts.

\subsection{Digital Twin Technology}

Digital twin technology creates virtual replicas of physical or logical systems that maintain synchronization with their real-world counterparts through continuous data exchange~\parencite{Shao_2021}. These virtual representations enable organizations to simulate system behavior, analyze potential changes, predict future states, and optimize performance without disrupting operational systems~\parencite{madni2018leveraging}. Digital twins combine real-time operational data with analytical models to provide dynamic, predictive capabilities that extend far beyond traditional monitoring and simulation approaches~\parencite{Khan_Saad_Niyato_Han_Hong_2022}.

As a concept, the digital twin originated in Product Lifecycle Management research. \citeauthor{Grieves_2023} traces the evolution of digital twins from their conceptual origins through contemporary applications, positioning digital twins as the integration of physical and virtual spaces that enables analysis, optimization, and prediction across system lifecycles~\parencite{Grieves_2023}. A defining distinction between digital twins and traditional simulation lies in synchronization. Where traditional simulations model hypothetical or designed system states, digital twins maintain continuous data exchange with their physical or logical counterparts. This exchange enables operational decision-making based on current rather than documented system states. Such synchronization directly addresses the documentation-reality gap that undermines Information Assurance and IT service management practices. Static models represent intended or designed system states; digital twins represent current operational states, updated continuously through integration with operational data sources. This distinction carries substantial implications for organizations that must make security and operational decisions based on accurate understanding of system configurations.

Standardization efforts have established frameworks for digital twin implementation that reduce interoperability concerns and provide structured guidance. The International Organization for Standardization published ISO 23247, the Digital Twin Framework for Manufacturing, defining a four-part reference architecture for observable elements, twin entities, data and device infrastructure, and cross-domain integration~\parencite{ISO_23247_2021}. \citeauthor{Shao_ISO_23247_2023} analyzed the standard in detail~\parencite{Shao_ISO_23247_2023}. While developed for manufacturing, the reference architecture establishes patterns of data synchronization, entity representation, and cross-domain integration that inform digital twin implementations in other domains. IEC 62832 complements ISO 23247 by addressing the industrial-process measurement, control, and automation aspects of digital twin frameworks~\parencite{Shao_2021}. In parallel, the Internet Engineering Task Force has published a draft reference architecture specifically addressing network infrastructure digital twins, representing direct engagement with the IT infrastructure domain~\parencite{IETF_NDT_2024}. IETF engagement signals that the networking community has recognized digital twin applicability for IT infrastructure, though the published architecture remains at draft stage and has not yet achieved the maturity of manufacturing-focused standards.

NIST has engaged with digital twin technology through NIST Internal Report 8356, finalized in February 2025, which addresses cybersecurity challenges and trust considerations for digital twin implementations~\parencite{NIST_IR_8356_2025}.
% Verified: NIST IR 8356 (Security and Trust Considerations for Digital Twin Technology) was published as final on February 14, 2025 (draft issued April 2021), per csrc.nist.gov/pubs/ir/8356/final, July 2, 2026. Bib entry year 2025 is correct.
Notably, this publication examines security considerations \textit{for} systems employing digital twins rather than digital twin applications \textit{to} Information Assurance, a distinction that reveals the current orientation of institutional research. NIST's Framework for Cyber-Physical Systems similarly addresses digital integration challenges at the intersection of physical and computational systems~\parencite{NIST_CPS_Framework_2017}. Together, these publications establish that authoritative institutions recognize digital twin relevance to security domains while treating digital twins primarily as systems requiring security rather than as tools for enhancing security operations. Reorienting from securing digital twins to employing them for security purposes represents a conceptual shift that the academic literature has begun but not yet completed.

Open source frameworks have emerged to reduce economic and technical barriers to digital twin adoption. The Digital Twin Consortium has established an open source collaboration initiative providing frameworks and reference implementations~\parencite{DTC_OpenSource_2024}, while the Eclipse Foundation offers Eclipse Ditto for IoT digital twins~\parencite{Eclipse_Ditto_2024} and Eclipse BaSyx for Asset Administration Shell industrial twins~\parencite{Eclipse_BaSyx_2024}. \citeauthor{Gil_2024} surveyed fourteen open source frameworks against ISO 23247 criteria and found substantial variation in maturity, documentation quality, and community support~\parencite{Gil_2024}. \citeauthor{Autiosalo_Siegel_Tammi_2021} introduced Twinbase, an open source server for the Digital Twin Web concept, enabling organizations to publish and discover digital twin descriptions through standardized interfaces~\parencite{Autiosalo_Siegel_Tammi_2021}. These open source options provide accessible entry points for organizations exploring digital twin capabilities, though their documentation and application examples focus on manufacturing and IoT rather than enterprise IT or Information Assurance contexts.

Digital twins support IT service delivery in ways that align with ITIL service design and service transition practices. Organizations can employ digital twins for capacity planning, change impact analysis, and service optimization: teams test changes and analyze performance implications before implementing modifications in production environments. Simulating proposed changes in a synchronized virtual environment reduces the risk of service disruptions while supporting more rapid and confident change implementation. When combined with digital thread traceability, digital twins enable organizations to assess the security and operational implications of proposed changes simultaneously. Current practices, which typically assess security and operational impacts through separate processes, cannot provide this integrated analysis capability.

Digital twins offer comparable advantages for Information Assurance, supporting security control validation, risk assessment, and incident response. Organizations can create digital twins of their information systems to simulate security scenarios, test control effectiveness, and analyze attack vectors in environments isolated from production operations. These virtual representations enable security teams to evaluate the impact of proposed security controls, assess the effectiveness of defensive measures, and predict system behavior under various threat scenarios. Digital twins support RMF assessment activities by enabling organizations to test security configurations and validate control implementations before deployment to production environments. Simulating adversary behavior against synchronized representations of actual system configurations enables security testing that reflects operational reality rather than documented assumptions.

\subsection{Product Lifecycle Management}

Product Lifecycle Management provides frameworks and toolsets for managing information, processes, and resources throughout a system's entire lifecycle from initial conception through retirement~\parencite{Ross_Winstead_McEvilley_2022}. PLM integrates data from diverse sources, maintains configuration baselines, manages change processes, and ensures that stakeholders access current, accurate information about system states and changes. This integrated approach to lifecycle management extends beyond simple version control to encompass configuration management, change coordination, release management, and information governance.

PLM originated in manufacturing and product development, where organizations confronted the challenge of managing complex products involving thousands of components, multiple suppliers, extended development timelines, and rigorous quality requirements. Manufacturing organizations developed PLM practices to address challenges that bear structural similarity to those confronting enterprise IT management: maintaining accurate documentation of complex configurations, coordinating changes across interdependent components, ensuring consistency between design documentation and as-built products, and managing compliance with regulatory requirements throughout product lifecycles. Over decades of industrial application, the commercial PLM ecosystem has matured, producing platforms that manage engineering artifacts across distributed supply chains.
% Softened: removed the unverifiable "large volumes" scale flourish; no authoritative public source quantifies artifact volume across commercial PLM platforms (analyst data such as CIMdata is paywalled), so the sentence now makes no quantitative claim requiring citation.
Open source PLM alternatives have emerged to reduce cost barriers, though \citeauthor{Laili_2024} noted that industrial open source PLM solutions vary considerably in maturity and require careful evaluation against organizational requirements~\parencite{Laili_2024}.

Applying PLM principles to information systems represents a conceptual extension from these manufacturing origins. Physical products have lifecycles that parallel information system lifecycles in important ways: conception, design, development, deployment, operation, maintenance, and retirement. PLM practices developed for managing physical product lifecycles address challenges of configuration management, change coordination, and baseline maintenance that information system managers confront daily. The question is whether PLM tools and methodologies can be adapted effectively for information system contexts, where system components are logical rather than physical, change cycles occur in hours rather than months, and deployment involves software configurations rather than manufacturing processes.

PLM capabilities align closely with ITIL service lifecycle management concepts. Organizations can implement PLM frameworks to support ITIL configuration management by maintaining authoritative configuration baselines and managing configuration item relationships. PLM change coordination capabilities enhance ITIL change management by providing improved visibility into change impacts across service dependencies. Maintaining integrated views of system configurations, security controls, and service delivery components reduces inconsistencies between security and IT operations teams, improves change coordination, and supports more effective compliance management across the system lifecycle. PLM's emphasis on managing information across organizational boundaries addresses the silo problem that fragments visibility in current IT environments: security teams, operations teams, and application teams can operate from a shared, authoritative information base rather than maintaining separate, potentially inconsistent documentation.

PLM principles apply equally within Information Assurance, addressing challenges in managing complex IT infrastructures and security postures throughout the system lifecycle. The RMF explicitly recognizes the importance of lifecycle management, requiring organizations to maintain security controls and documentation throughout system operation and into decommissioning. PLM approaches support these requirements by managing security control baselines, coordinating changes across interconnected systems, maintaining configuration integrity, and ensuring that security teams operate from consistent, current information throughout the authorization boundary. PLM configuration baseline capabilities directly address the requirement in NIST SP 800-53 control CM-2 for documented, maintained system baselines~\parencite{Force_2020}. PLM change coordination capabilities support control CM-3 requirements for configuration change control. PLM access governance capabilities align with CM-5 requirements for access restrictions governing system changes. These correspondences between PLM capabilities and specific compliance requirements suggest that PLM practices could address compliance obligations that organizations currently struggle to fulfill through manual processes.

\subsection{Institutional Endorsement and Strategic Direction}

That Digital Engineering has matured from experimental practice to institutional mandate provides important context for evaluating its potential applicability beyond defense and aerospace. The breadth and depth of institutional endorsement establish Digital Engineering as a mature discipline validated by the organizations most experienced in managing complex, compliance-intensive systems. Understanding this institutional foundation informs assessment of whether Digital Engineering merits investigation for enterprise IT and Information Assurance applications.

The Digital Engineering Strategy, published in June 2018 under the institution's then-current name of Department of Defense, established the formal vision for transforming defense acquisition through Digital Engineering practices~\parencite{DoD_DE_Strategy_2018}. Five strategic goals collectively articulate the transformation Digital Engineering envisions: formalize the development, integration, and use of models to inform enterprise and program decision-making; provide an authoritative source of truth by establishing a digital thread as the means of delivering the technical baseline; incorporate technological innovation to improve the engineering practice; establish a Digital Engineering ecosystem to provide the infrastructure and environments for performing Digital Engineering activities; and transform the culture and workforce to adopt and support Digital Engineering across the lifecycle. These goals extend beyond technology adoption to encompass organizational transformation, workforce development, and cultural change. Any transfer of Digital Engineering to enterprise IT contexts must address these same dimensions.

DoDI 5000.97, issued in December 2023, codified these strategic aspirations into mandatory requirements for defense programs~\parencite{DoDI_5000_97_2023}. It mandates that programs use digital artifacts as the authoritative source of system information and maintain digital thread capabilities throughout the acquisition lifecycle. Moving from strategic vision to regulatory mandate marks a major institutional commitment: defense programs must now implement Digital Engineering practices rather than merely consider them. Complementing this mandate, the Systems Engineering Guidebook provides detailed implementation guidance for meeting these requirements, establishing processes, methods, and tool expectations for Digital Engineering practice within defense acquisition~\parencite{DoD_Engineering_2022}. A clear progression from strategy document to formal instruction to implementation guidebook demonstrates institutional maturation that parallels the compliance framework progression familiar in Information Assurance contexts, which moves from policy through standards through implementation guidance.

NASA has pursued a parallel path toward Digital Engineering adoption, lending independent validation to the approach. NASA-HDBK-1004, the NASA Digital Engineering Acquisition Framework Handbook, establishes guidance for incorporating Digital Engineering practices into NASA programs~\parencite{NASA_HDBK_1004_2020}. A companion document, the NASA Future Model-Based Systems Engineering Vision and Strategy Bridge, articulates NASA's long-term vision for MBSE adoption across the agency's portfolio of complex, mission-critical programs~\parencite{NASA_MBSE_Vision_2021}. That NASA independently developed Digital Engineering guidance carries particular significance. Both the Department of War and NASA confronted similar challenges: complex systems demanding comprehensive documentation, stringent compliance requirements, and mission-critical operations tolerating no ambiguity in system understanding. Both organizations converged on Digital Engineering as the solution. This independent convergence strengthens the argument that Digital Engineering addresses fundamental challenges in managing complex, compliance-intensive systems rather than representing a domain-specific solution applicable only to defense acquisition.

INCOSE has positioned Digital Engineering as the future direction of the systems engineering discipline. Its Systems Engineering Vision 2035 establishes a roadmap for transforming systems engineering practice through Digital Engineering, model-based methods, and integrated digital environments~\parencite{INCOSE_Vision_2035}. Complementing this vision, the INCOSE Systems Engineering Handbook, Fifth Edition, provides comprehensive guidance for implementing systems engineering practices including MBSE, reflecting the discipline's current state of practice and serving as the primary professional reference for systems engineers worldwide~\parencite{INCOSE_Handbook_2023}. INCOSE's Digital Engineering Information Exchange Working Group specifically addresses the interoperability and integration challenges that arise when organizations attempt to exchange Digital Engineering artifacts across organizational and tool boundaries~\parencite{DigitalEngineeringInformationExchange}. This working group's existence acknowledges that Digital Engineering implementation requires addressing practical interoperability challenges, and those challenges would intensify when extending Digital Engineering beyond its established systems engineering context into enterprise IT domains.

The Systems Engineering Body of Knowledge (SEBoK) is jointly managed by INCOSE, the IEEE Systems Council, and the Stevens Institute Systems Engineering Research Center. It provides the globally recognized authoritative reference defining Digital Engineering's relationship to MBSE, digital threads, and authoritative sources of truth within the ISO/IEC/IEEE 15288:2023 systems engineering lifecycle framework~\parencite{SEBoK_2024}. SEBoK establishes the conceptual vocabulary and architectural relationships that the systems engineering community employs when discussing Digital Engineering, providing the definitional foundation on which this dissertation's investigation builds.

SERC has produced a body of research publications supporting Digital Engineering implementation and assessment. Its Digital Engineering Competency Framework defines 962 Knowledge, Skills, Abilities, and Behaviors organized by proficiency levels, providing the most comprehensive specification of what professionals require to practice Digital Engineering effectively~\parencite{Hutchinson_2021}. A companion Digital Engineering Metrics framework establishes measurement approaches for assessing adoption progress and benefits realization, addressing the measurement challenge that the literature identifies as a persistent barrier to evidence-based adoption decisions~\parencite{Hutchison_2020}. Additionally, the SERC Systems Engineering Modernization report examines the integration of Digital Engineering with Mission Engineering, Agile, and DevOps practices, establishing connections between Digital Engineering and the development methodologies prevalent in enterprise IT contexts~\parencite{SERC_SE_Modernization_2022}. Collectively, these SERC publications provide a research infrastructure of competency definitions, metrics, and integration frameworks that would support Digital Engineering adoption in new domains, including enterprise IT and Information Assurance.

NIST's contributions span multiple relevant domains. The Framework for Cyber-Physical Systems addresses integration challenges at the intersection of physical and computational systems~\parencite{NIST_CPS_Framework_2017}. NIST's systems security engineering publications reference model-based approaches and traceability requirements that align with Digital Engineering principles. These publications comprise Special Publication 800-160 Volume 1 Revision 1 on Engineering Trustworthy Secure Systems~\parencite{Ross_Winstead_McEvilley_2022} and Volume 2 Revision 1 on Developing Cyber-Resilient Systems~\parencite{NIST_SP_800_160_V2_2021}. NIST Internal Report 8356 on digital twin security considerations~\parencite{NIST_IR_8356_2025} and the Open Security Controls Assessment Language initiative~\parencite{NIST_OSCAL_2023} further demonstrate NIST engagement with concepts that Digital Engineering could integrate. Taken together, these publications establish that the institution responsible for federal cybersecurity standards recognizes the relevance of model-based approaches, traceability, and digital integration to security engineering. NIST has not yet explicitly connected these concepts under the Digital Engineering framework, however, nor provided specific guidance for their application in enterprise IT contexts. This dissertation addresses one such research gap: the distance between NIST's recognition of relevant principles and the absence of specific Digital Engineering implementation guidance for enterprise IT.

\subsection{Enterprise Architecture Frameworks}

Enterprise architecture frameworks provide the structural foundation on which Digital Engineering implementations rest. These frameworks define how organizations represent, analyze, and manage the complex relationships among business processes, information flows, application systems, and technology infrastructure. Understanding the enterprise architecture framework ecosystem is essential for evaluating Digital Engineering's potential in enterprise IT contexts because the frameworks determine what can be modeled, how models are organized, and what analytical capabilities models support.

The Unified Architecture Framework, now codified as ISO/IEC 19540 through the Object Management Group~\parencite{OMG_UAF_2022}, is a major consolidation of enterprise architecture standards. UAF merges three military architecture frameworks: the Department of Defense Architecture Framework (DoDAF)~\parencite{DoDAF_2009}, the UK Ministry of Defence Architecture Framework (MODAF), and the NATO Architecture Framework~\parencite{NATO_NAF_2018}. This consolidation proceeded from the explicit recognition that the architectural concepts developed for military system-of-systems management possess broad applicability beyond defense contexts. The Object Management Group determined that approximately 90\% of concepts captured in military architecture frameworks prove equally applicable in commercial domains~\parencite{OMG_UAF_About_2024}, motivating the development of a unified framework that serves both defense and commercial organizations without requiring domain-specific adaptations.

UAF employs a grid-based structure defining seventy-one view specifications organized through two complementary parts: the UAF Domain Metamodel (DMM), which defines the underlying concepts, relationships, and constraints; and the UAF Modeling Language (ML), which provides the concrete notation for expressing architecture descriptions~\parencite{OMG_UAF_Spec_2022}. Grid rows represent stakeholder domains, namely Strategic, Operational, Services, Personnel, and Resources, while columns represent architecture aspects including taxonomy, structure, connectivity, processes, states, interaction, constraints, and traceability. Together, this grid organization enables architects to represent any complex system-of-systems from multiple stakeholder perspectives simultaneously, with formal relationships connecting views across the grid to maintain consistency. The traceability column explicitly addresses the requirement for maintaining verified connections across architectural perspectives, a requirement that aligns directly with the digital thread concept and with the compliance traceability obligations that NIST SP 800-53 imposes.

Evolution from domain-specific frameworks to the unified standard proceeded through several stages. DoDAF 2.0 established the foundational metamodel for defense architecture description~\parencite{DoDAF_2009}, MODAF developed parallel conventions for UK defense, and the NATO Architecture Framework provided capabilities for alliance-wide interoperability~\parencite{NATO_NAF_2018}. Earlier evaluation identified limitations in these individual frameworks: the National Defense Industrial Association found that DoDAF required augmentation to support systems engineering activities~\parencite{NDIA_SE_DoDAF_2011}, and \citeauthor{Hause_2010} identified specific gaps between DoDAF's metamodel and Model-Based Systems Engineering practice~\parencite{Hause_2010}. The Unified Profile for DoDAF and MODAF (UPDM) was an initial consolidation effort that evolved into UAF under broader OMG stewardship, with ISO/IEC standardization following to establish UAF as a globally recognized standard~\parencite{OMG_UAF_UPDM_2017}.

Comparative research validates UAF's position relative to alternative enterprise architecture frameworks. \citeauthor{Bankauskaite_2019} evaluated enterprise architecture frameworks using weighted criteria encompassing completeness, maturity, standardization, and tool support, finding that UAF achieved the highest overall rating among the frameworks evaluated, surpassing TOGAF~\parencite{TOGAF_2018}, DoDAF, MODAF, NAF, and FEAF~\parencite{Bankauskaite_2019}. This comparative advantage derives from UAF's comprehensive metamodel, its ISO international standardization, and its explicit design for both defense and commercial applications. Research by \citeauthor{Eichmann_Melzer_God_2019} demonstrated UAF application to system-of-systems development~\parencite{Eichmann_Melzer_God_2019}. In a military context, \citeauthor{Liu_2023} applied UAF to system-of-systems design using top-down MBSE methodology~\parencite{Liu_2023}. Building on UAF-based MBSE, \citeauthor{Torkjazi_2024} integrated autonomy into system-of-systems architectures~\parencite{Torkjazi_2024}, while \citeauthor{Abhaya_2021} established methods for building system-of-systems models~\parencite{Abhaya_2021}. Collectively, these studies demonstrate UAF's growing adoption and validation across diverse system-of-systems contexts, though applications to enterprise IT infrastructure and Information Assurance remain absent from the literature identified in this review.

The complementary relationship between UAF and the commercially prevalent Open Group Architecture Framework (TOGAF) reduces adoption barriers for organizations already invested in enterprise architecture practice. TOGAF provides a comprehensive architecture development method and enterprise architecture governance framework widely adopted in commercial IT contexts~\parencite{TOGAF_2018}. A joint white paper between The Open Group and MITRE Corporation established that UAF and TOGAF address different but complementary aspects of enterprise architecture: TOGAF provides the methodology and governance framework for architecture development, while UAF provides the architecture description specification for expressing what has been architected~\parencite{OpenGroup_MITRE_2013}. Organizations can retain TOGAF methodology while adopting UAF for architecture description, enabling incremental adoption rather than wholesale replacement of established practices. This compatibility matters especially for enterprise IT organizations that have invested in TOGAF-based architecture programs: it suggests that UAF adoption for model-based architecture description need not require abandonment of existing architecture governance investments.

UAF's foundation on UML and SysML profiles enables model-based documentation approaches that address the accuracy and currency challenges plaguing traditional enterprise architecture implementations~\parencite{Aerospace_UAF_2023}. Because UAF architecture descriptions are expressed as formal models rather than static documents, they inherit the automated consistency checking, impact analysis, and traceability capabilities that model-based approaches provide. A model-based foundation connects enterprise architecture practice to the broader Digital Engineering ecosystem: UAF models can participate in digital threads, maintain synchronization with operational systems through digital twin mechanisms, and be managed through PLM lifecycle processes. UAF has been adopted by the Department of War~\parencite{DoD_Engineering_2022}, by NATO~\parencite{NATO_NAF_2018}, and by the UK Ministry of Defence, and it carries ISO international standardization as well. Together these endorsements establish UAF as the authoritative consolidating framework for organizations that need architecture capabilities across multiple domains. For enterprise IT and Information Assurance applications, UAF provides structured approaches to documenting systems, relationships, and security requirements that align with both NIST and ITIL expectations. Understanding this enterprise architecture foundation provides essential context for evaluating how Digital Engineering extends architectural approaches from static documentation toward living, model-based representations that maintain currency with operational reality.

\section{Gaps in Current Practice}

Despite advances in Information Assurance methodologies and IT service management frameworks, organizations continue to encounter challenges that limit their effectiveness in protecting information assets and delivering reliable services. The NIST Risk Management Framework and ITIL provide structured approaches to Information Assurance and IT service management respectively, yet implementation challenges persist across both domains. Digital Engineering practices, as the 2018 Digital Engineering Strategy directs for defense programs, are designed to address structurally analogous challenges: maintaining authoritative documentation, ensuring configuration visibility, and providing traceability across complex systems~\parencite{DoD_DE_Strategy_2018}. Whether these capabilities translate to enterprise IT and Information Assurance contexts, and whether professionals in those domains recognize the potential relevance, remains uninvestigated. Examining the specific gaps that persist in current practice illuminates the structural parallels that motivate this investigation.

These gaps persist despite framework maturity and organizational investment, echoing the pattern documented earlier in this chapter: these outcomes are consistent with limitations inherent in document-centric methodologies rather than with execution failures alone.

\subsection{Information Assurance Challenges}

Organizations implementing the NIST Risk Management Framework struggle to maintain visibility into their security postures across complex, distributed information systems. The RMF continuous monitoring phase requires organizations to maintain ongoing awareness of security control effectiveness and system security state. Security teams often lack accurate, current understanding of system configurations, security control implementations, and the dependencies that affect security effectiveness. Visibility gaps manifest throughout the RMF lifecycle: organizations find it difficult to track security dependencies effectively, leaving vulnerabilities unidentified when teams implement changes. When teams cannot maintain accurate documentation of system interconnections and data flows, particularly in environments with rapid deployment cycles, the risk assessment and incident response capabilities that the RMF demands break down.

Understanding cascading impacts poses a further challenge. When incidents occur or defenders discover vulnerabilities, they must rapidly assess which systems are affected, what data are at risk, and how compromise of one system might enable access to interconnected systems. Such assessment requires understanding of system dependencies that current documentation approaches cannot maintain. When defenders cannot trace first, second, and third order impacts, incident response becomes a broad search effort rather than a precision operation, consuming time and resources while adversaries retain the initiative.

Security control implementation presents distinct challenges in modern enterprise environments characterized by hybrid cloud deployments, distributed architectures, and frequent changes. Organizations must implement and maintain consistent security controls from NIST SP 800-53 across diverse platforms while supporting continuous deployment practices and rapid update cycles~\parencite{Force_2020}. Traditional security configuration management approaches fail to scale effectively in these dynamic environments, leading to inconsistent security postures and compliance gaps. The challenge compounds when organizations attempt to validate control effectiveness across interconnected systems where authorization boundaries grow increasingly complex to define and maintain.

\subsection{IT Service Management Challenges}

IT service management faces parallel challenges in maintaining accurate system documentation. As noted earlier in this chapter, most CMDB implementations fail to deliver intended value, leaving organizations without authoritative sources for configuration information~\parencite{Gartner_CMDB_2019}. Manual configuration tracking cannot keep pace with the rate of change in modern IT environments. Shadow IT creates blind spots where undocumented systems introduce unknown dependencies and security risks. Change management processes suffer when impact assessments rely on incomplete or inaccurate dependency information.

The economic dimensions of these failures merit examination. Organizations commit budget, staff time, and tooling to CMDB implementations, documentation efforts, and change management processes.
% Softened: removed the "considerable resources" magnitude claim; no citable public analyst figure quantifies organizational CMDB and documentation investment (Gartner CMDB research is paywalled), so the sentence now states only the uncontroversial fact that such investments occur, which the already-cited Gartner_CMDB_2019 failure statistic presupposes.
When these investments fail to deliver intended value, organizations face difficult choices: invest additional resources attempting to improve failing approaches, accept degraded capabilities and increased risk, or seek alternative approaches that address root causes rather than symptoms. Digital Engineering is a candidate alternative approach whose reported benefits in defense and aerospace contexts suggest potential applicability to enterprise IT environments.

The convergence of these challenges creates a compounding effect where neither Information Assurance nor IT service management can achieve their objectives independently. Security teams cannot effectively assess risks without accurate understanding of IT infrastructure. IT teams cannot effectively manage changes without understanding security implications. Both domains require the visibility and accurate documentation that current practices fail to provide.

\subsection{The Documentation-Reality Gap}

The persistent gap between documentation and operational reality, documented in the Challenges in Current Practice discussion earlier in this chapter, forms the common thread connecting failures across both domains. When documentation diverges from reality, security assessments measure fiction rather than fact, change impact analyses miss undocumented dependencies, incident responders lose time reconciling records with operational systems, and compliance auditors cannot verify that documented controls exist in practice. Digital Engineering offers capabilities designed to address this gap through authoritative sources of truth, continuous synchronization between models and operational systems, and automated consistency verification. The 2018 Digital Engineering Strategy directs defense programs to establish these capabilities in complex, compliance-intensive environments~\parencite{DoD_DE_Strategy_2018}. The question this research investigates is whether IT and Information Assurance professionals recognize the potential value of these capabilities for their work, a prerequisite for determining whether Digital Engineering practices developed in defense contexts merit investigation for broader enterprise application.

\section{Research Questions}

Two findings now stand in tension. Defense and aerospace have committed to Digital Engineering at institutional scale to manage the very complexity that overwhelms enterprise IT and Information Assurance, and the limited measured evidence available is encouraging. Yet no study has measured whether the professionals who manage that complexity recognize Digital Engineering as a potential remedy. That unmeasured question defines this study's focus. Building on the challenges documented in current practice and the capabilities offered by Digital Engineering, this investigation pursues three questions:

\begin{enumerate}
  \item To what extent are Information Technology and Information Assurance professionals aware of Digital Engineering capabilities, including Model-Based Systems Engineering, digital threads, digital twin technologies, and Product Lifecycle Management principles?

  \item To what extent do Information Technology and Information Assurance professionals perceive Digital Engineering capabilities as potentially valuable or important for their work in Information Assurance, security compliance, and IT service delivery?

  \item To what extent do Information Technology and Information Assurance professionals believe that Digital Engineering practices could help them in performing their jobs, meeting compliance requirements, or enhancing organizational capabilities in Information Assurance and IT service delivery?
\end{enumerate}

Research Question 3 is examined through the operational and compliance-oriented items distributed across Sections~3 through~5 of the survey instrument, grouped into IT-operational, security-operational, and explicitly compliance-framed sets, so that endorsement of Digital Engineering's job-performance and capability value can be read separately from endorsement of its compliance value.

These questions are posed of the profession at large, and the study answers them for the population its recruitment actually reached: a non-probability sample drawn through defense- and aerospace-weighted professional channels and weighted toward earlier-career respondents. Chapter~3 specifies that frame and Chapter~5 bounds the conclusions to it.

\subsection{Awareness and Perception as the Necessary First Step}

These research questions focus on professional awareness and perceptions as the groundwork for subsequent inquiry. Establishing awareness levels and perceived value is an essential first step. It must precede any investigation of practical implementation approaches, organizational adoption strategies, or empirical validation of Digital Engineering benefits in Information Assurance and IT service management contexts. Preliminary evidence from adjacent domains suggests that such investigation is urgently needed: \citeauthor{Henderson_McDermott_Salado_2024} found that approximately 22\% of systems engineering professionals, all practitioners within the discipline that developed MBSE, could not articulate a clear definition of MBSE~\parencite{Henderson_McDermott_Salado_2024}. If awareness deficits persist within the originating discipline, awareness among professionals in Information Assurance and IT Service Management, disciplines that have never systematically encountered these methodologies, requires empirical measurement rather than assumption.

\subsection{Distinguishing General Value from Personal Applicability}

The distinction between Research Question 2 and Research Question 3 requires clarification, because both address professional perceptions of Digital Engineering value. Research Question 2 measures whether professionals recognize Digital Engineering capabilities as abstractly valuable or important for their professional domains, an assessment of general relevance that does not require the respondent to evaluate specific operational impact within their own work context. Research Question 3 probes a more personal and practical construct: whether professionals believe Digital Engineering could tangibly help them perform their jobs, meet the compliance requirements they personally face, or enhance the specific organizational capabilities they are responsible for delivering. Technology adoption literature distinguishes between recognizing that a capability has general value and believing that the same capability would improve one's own work. Professionals may acknowledge that Digital Engineering offers theoretical value for their domain while remaining skeptical about its practical applicability to their specific organizational context. This distinction carries substantial implications for adoption strategy: a professional community that perceives general value but doubts personal applicability requires different intervention than one that perceives both.

Whether the two are empirically separable in the responses, or instead express a single perceived-value appraisal at different levels of specificity, is itself an empirical question. The analysis does not construct separate composites for general and personal value. Chapter~4's construct-validity results collapse the IT-value and security-value dimensions into a single perceived-value construct. That outcome indicates the two questions address one perceived-value appraisal at different levels of specificity, and Chapters~4 and~5 report them on that basis.

\section{Theoretical Framework and Research Propositions}

The investigation draws on three established streams of technology-adoption theory to connect its constructs. The Technology Acceptance Model holds that behavioral intention to adopt a technology is driven by perceived usefulness and perceived ease of use~\parencite{Davis_1989}. The Unified Theory of Acceptance and Use of Technology generalizes this into performance expectancy, effort expectancy, social influence, and facilitating conditions, with individual characteristics acting as moderators~\parencite{Venkatesh_2003}. Rogers' Diffusion of Innovations situates adoption within the perceived attributes of an innovation---among them relative advantage, compatibility, and observability---and emphasizes that perceptions rather than objective properties govern uptake~\parencite{Rogers_2003}. \citeauthor{Call_Herber_2022} mapped this framework directly onto Model-Based Systems Engineering adoption, establishing its applicability to the present domain~\parencite{Call_Herber_2022}.

These streams motivate the conceptual model in Figure~\ref{fig:conceptual}. The model positions awareness of the discipline and recognition of its capabilities, the focus of Research Question~1, as antecedents of perceived value, the focus of Research Questions~2 and~3. This construct corresponds to perceived usefulness in the Technology Acceptance Model and to relative advantage in the diffusion framework. Perceived value in turn informs adoption intent, captured here through interest in learning more and willingness to recommend organizational exploration. Professional field and experience enter the model as moderators, consistent with the individual-difference moderators of the Unified Theory, while self-reported familiarity is treated as a covariate of perceived value. The survey instrument operationalizes the perceived-value slice of these frameworks together with two coarse indicators of adoption intent. It does not measure effort expectancy, facilitating conditions, or social influence, so the study tests the value-perception portion of the adoption models rather than the models in full; the remaining constructs are left to the future work outlined in Chapter~5.

\begin{figure}[ht]
  \centering
    \includegraphics[width=0.92\linewidth, alt={Path diagram of the study's conceptual model, running left to right through three boxes. Awareness and capability recognition, labeled Research Question 1, connects by an arrow annotated TAM and UTAUT to perceived value, labeled Research Questions 2 and 3, which connects by an arrow annotated Diffusion of Innovations to adoption intent. A fourth box below, holding professional field and experience, sends dashed arrows to both paths to mark them as moderators.}]{figures/fig_conceptual_model.pdf}
  \caption{Conceptual model linking awareness and capability recognition to perceived value and adoption intent, with professional field and experience as moderators.}
  \label{fig:conceptual}
\end{figure}

Because the design is cross-sectional and the application domain is largely unstudied, the research advances propositions rather than formal causal hypotheses. Four directional expectations follow from the framework and the gaps documented above:

\begin{enumerate}
  \item \textbf{P1 (Low awareness).} Awareness of and recent exposure to Digital Engineering are low among IT and Information Assurance professionals, who sit outside the systems engineering discipline in which these methods developed.
  \item \textbf{P2 (Value despite low awareness).} Perceived value is nonetheless positive once capabilities are described in functional terms, because Digital Engineering addresses the documentation, traceability, and visibility problems these professionals routinely face. Because an instrument that measures perceived value must necessarily describe the capability being evaluated, this proposition carries a demand-characteristics risk; the neutral-framing safeguards and bias diagnostics detailed in Chapters~3 and~4 are designed to test whether the resulting endorsement survives that scrutiny.
  \item \textbf{P3 (Variation by field and experience).} Perceived value and awareness vary systematically with professional field and experience level rather than being uniform across the community.
  \item \textbf{P4 (Familiarity and value).} Greater familiarity with Digital Engineering is associated with higher perceived value.
\end{enumerate}

Chapter~4 evaluates each proposition against the survey data, and Chapter~5 returns to them in stating the conclusions. Propositions~P1 and~P2 correspond primarily to Research Questions~1 and~2, with P2's expectation of functional-value endorsement extending to the applicability beliefs of Research Question~3; P3 and P4 cut across all three questions as expected patterns of variation and association.

\section{Research Scope and Approach}

This research examines professional perceptions in four areas: Model-Based Systems Engineering, digital threads, digital twin technology, and Product Lifecycle Management. The investigation assesses awareness and perceived value of Model-Based Systems Engineering approaches for representing information system architectures and security controls. It asks to what extent professionals perceive value in digital threads for maintaining authoritative traceability between security requirements, control implementations, and compliance evidence as required by frameworks such as the NIST RMF\@. It explores perceptions of digital twin capabilities for security simulation, testing, and IT service modeling. Finally, it investigates how strongly professionals value Product Lifecycle Management principles for managing information system configurations and security control baselines throughout the system lifecycle.

\subsection{Methodological Approach}

The research employs a quantitative survey method to collect data from IT and Information Assurance professionals. Survey methodology enables systematic data collection from a broad population of practitioners, supporting statistical analysis and generalization of findings. The anonymous nature of survey research encourages candid responses about professional knowledge gaps and organizational capabilities. Chapter 3 presents the complete research methodology including survey design, sampling strategy, and analytical approach.

\subsection{Target Population}

The research targets IT and Information Assurance professionals across diverse organizational contexts. This population includes security analysts, compliance officers, IT managers, network administrators, systems administrators, and other professionals who implement and maintain Information Assurance controls or deliver IT services. By surveying professionals across organizational contexts rather than limiting investigation to a single industry or organization, the research enables assessment of awareness and perceived value across the professional community.

This target population intentionally extends beyond the systems engineering professionals who have been the focus of previous MBSE adoption research. While studies by \citeauthor{Call_2024} and \citeauthor{Henderson_McDermott_Salado_2024} examined perceptions among systems engineering professionals~\parencite{Call_2024, Henderson_McDermott_Salado_2024}, no comparable research addresses IT and Information Assurance practitioners. This population is a crucial test of Digital Engineering's transferability: if professionals who have never encountered these methodologies nevertheless recognize their potential value for addressing familiar challenges, the case for cross-disciplinary investigation strengthens substantially.

\subsection{Potential Benefits for Under-Resourced Organizations}

The potential transferability of Digital Engineering benefits carries particular significance for small and under-resourced organizations. Many under-resourced organizations require effective Information Assurance and IT service delivery capabilities. These include Federally Qualified Health Centers and safety-net hospitals serving low-income patients, tribal health and government organizations, rural healthcare providers, school districts and minority-serving institutions in under-resourced communities, public libraries, small and rural municipalities, and the non-profit and social-service entities that address community needs. For these organizations the frameworks already exist; what most sharply constrains them is the scarcity of the specialized labor those frameworks presuppose. Model-Based Systems Engineering has already been applied to rural-healthcare planning in exactly such resource-limited settings, which establishes a concrete precedent for carrying these methods beyond well-funded enterprises~\parencite{Berg_2023}.

Under-resourced organizations must frequently demonstrate compliance with regulatory frameworks, security standards, and funding requirements. Healthcare providers must satisfy HIPAA security requirements. Educational institutions must protect student data under FERPA\@. Social service organizations must safeguard client information while demonstrating accountability to funding agencies. These compliance obligations demand documented evidence of security controls and system configurations, and producing that documentation consumes scarce staff time and organizational resources.

If Digital Engineering practices can reduce the burden of compliance documentation while improving documentation accuracy, organizations with limited resources could redirect staff effort toward direct service provision rather than documentation administration. Digital threads, through automated traceability, could reduce the manual effort required for compliance audits. MBSE could maintain documentation accuracy through automated synchronization rather than manual updates. PLM configuration management capabilities could, in turn, lower the specialized expertise required to maintain accurate system documentation. Each of these mechanisms substitutes reusable, automatable model structure for scarce specialized labor. To the extent that the substitution holds, it could place authoritative documentation and demonstrable compliance within reach of organizations that cannot today afford the dedicated staff such work requires; whether that capability also narrows the broader equity gap is the distributional question taken up under Significance.

This research does not presume that Digital Engineering benefits will transfer effectively to resource-constrained organizations. Accordingly, the survey population targets IT and Information Assurance professionals broadly, not exclusively those working in under-resourced settings. However, by establishing baseline awareness and perceived value data across the professional community, this research creates a foundation for subsequent investigation of Digital Engineering applicability in diverse organizational contexts. If professionals perceive value in Digital Engineering capabilities, future research can examine practical implementation approaches suitable for organizations with varying resource levels.

The logical pathway from defense and aerospace demonstration to broader application proceeds through several steps. First, defense and aerospace organizations demonstrate that Digital Engineering delivers measurable benefits for complex systems with stringent compliance requirements. Second, research establishes whether IT and Information Assurance professionals outside these sectors recognize potential value in Digital Engineering capabilities. Third, if perceived value exists, subsequent research can examine implementation approaches, adaptation requirements, and cost-benefit considerations for organizations in different contexts. This research addresses the second step: determining whether professional awareness and perceived value support continued investigation of Digital Engineering for enterprise IT and Information Assurance applications.

\subsection{Why Perceptions Matter}

Investigating professional perceptions requires justification, given the availability of alternative research approaches. Technology adoption research consistently shows that perceived value influences adoption decisions regardless of actual value~\parencite{Davis_1989}. Professionals who do not perceive value in a capability will not advocate for its adoption within their organizations. Establishing whether IT and Information Assurance professionals recognize potential value in Digital Engineering capabilities therefore addresses a prerequisite question for successful adoption.

Furthermore, perception research enables assessment of awareness gaps that might impede adoption. If professionals are unaware of Digital Engineering capabilities, education and communication initiatives become necessary precursors to adoption efforts. If professionals are aware but do not perceive value, the theoretical premise that Digital Engineering addresses recognized needs requires reconsideration. Understanding the current state of professional awareness and perceptions enables targeted strategies for advancing Digital Engineering adoption in Information Assurance and IT service management domains.

\subsection{Relationship to Prior Work}

This research builds on a small body of prior work while extending investigation to a professional population that earlier studies did not reach. By surveying professionals actively working in these domains, it identifies whether practitioners recognize connections between their current practices and Digital Engineering capabilities. Findings illuminate whether existing Information Assurance and IT service delivery frameworks already incorporate concepts analogous to Model-Based Systems Engineering, digital threads, digital twins, or Product Lifecycle Management, or whether these Digital Engineering capabilities are genuinely novel approaches within Information Technology contexts. Such understanding establishes whether Digital Engineering offers new conceptual frameworks for addressing Information Assurance and IT service delivery challenges or whether it primarily provides different terminology for existing practices.

\citeauthor{Bonar_Hastings_2024} established an initial reference model showing how Digital Engineering practices can support and enhance compliance verification within the context of information systems~\parencite{Bonar_Hastings_2024}. The current research extends this foundation by examining whether the broader professional community recognizes value in the capabilities that the reference model proposes.

\section{Significance of the Research}

This research carries implications across academic, industry, national, and societal dimensions. Understanding these dimensions of significance contextualizes the contribution this investigation makes to knowledge and practice.

\subsection{Academic Significance}

The academic significance of this research begins with an absence. As Chapter 2 documents, the literature review identified a near-complete absence of academic research applying Model-Based Systems Engineering, digital threads, digital twins, or Product Lifecycle Management to enterprise IT infrastructure or Information Assurance programs. This study addresses that absence directly by establishing whether the professionals who work in these domains recognize a gap that Digital Engineering may address. This gap persists even though NIST and ITIL frameworks impose requirements for capabilities that Digital Engineering could provide: enterprise architecture representation, documentation accuracy, and traceability.

This research contributes empirical evidence to a domain currently characterized by theoretical proposition rather than data. Perception data collected from working professionals enables assessment of whether the theoretical framework connecting Digital Engineering capabilities to documented IT and Information Assurance challenges resonates with practitioners who experience those challenges directly. Findings indicating that professionals recognize value in Digital Engineering capabilities would validate the theoretical premise and support subsequent implementation research. Findings indicating limited awareness or skepticism would redirect the research agenda toward education initiatives or reconsideration of transferability assumptions. Either outcome advances the academic understanding of Digital Engineering's potential applicability beyond its established defense and aerospace foundations. The layered awareness-and-recognition design and the anonymous open-link mode of data collection also yield methodological results of independent interest. Foremost among them, as Chapter~4 details, is a validated and openly released response-authenticity screen for detecting coordinated and automated responding in anonymous online surveys, offered as a reusable methodological contribution in its own right.

\subsection{Industry Significance}

For industry, an adoption effort would meet either a receptive professional community or an indifferent one, and at present no data indicate which. This research supplies that missing signal by measuring how practitioners perceive Digital Engineering capabilities. Organizations considering Digital Engineering adoption can benefit from understanding current awareness levels and perceived value within the professional community. If research reveals widespread recognition of Digital Engineering value, organizations may find receptive audiences for adoption initiatives. If research reveals limited awareness or skepticism, organizations can anticipate the education and change management challenges that adoption would require.

Beyond individual organizations, the research characterizes which segments of the professional community perceive the least value and remain least aware, and therefore where receptiveness is weakest. Such information enables tool vendors, service providers, and standards organizations to focus communication and enablement on the audiences that the data identify as least reached, rather than allocating effort uniformly. Understanding professional perceptions enables more effective resource allocation across the ecosystem supporting Digital Engineering adoption.

\subsection{National Significance}

At the national scale, the documentation and visibility failures examined here cease to be an operational inconvenience and become a matter of national security and the protection of critical infrastructure. Federal information systems and national security systems protect information assets and enable government operations on which citizens depend. Organizations operating these systems face the challenges documented throughout this dissertation: maintaining accurate documentation, implementing consistent security controls, and verifying compliance across complex technical environments.

The 2018 Digital Engineering Strategy directs defense programs to establish digital threads, which are designed to enhance compliance verification and security assurance by providing authoritative traceability: verified connections between security requirements, control implementations, and compliance evidence~\parencite{DoD_DE_Strategy_2018}. Operators of federal and national security systems must demonstrate compliance with many regulatory frameworks and security standards, often requiring extensive manual effort to collect evidence and prepare for audits. Digital Engineering practices offer the potential to reduce the burden of compliance verification while improving the accuracy and currency of compliance documentation, enabling organizations to redirect resources toward proactive security improvements rather than documentation effort.

Understanding first, second, and third order impacts of security incidents is especially important for critical infrastructure protection. When adversaries compromise systems supporting government functions or critical infrastructure, defenders must rapidly assess the scope of compromise and potential cascading effects. Digital Engineering practices are designed to provide the visibility and traceability that enable rapid, accurate impact assessment, capabilities that current approaches lack. The present study measures professional receptiveness to this class of traceability capability rather than to cascading-impact assessment as a separately worded use case, a distinction Chapter~5 makes explicit. Whether organizations outside defense and aerospace can realize these capabilities effectively remains an open question that this research begins to address through professional perception data.

\subsection{Societal Significance}

Beyond organizations operating national security systems, Digital Engineering capabilities may benefit organizations serving communities with limited resources. Healthcare providers, educational institutions, social service organizations, and other under-resourced entities face the same Information Assurance and IT service delivery challenges as large enterprises, often with fewer resources to address them.

The distributional character of the burden gives the issue its equity significance. The documentation, traceability, and verification labor that compliance demands falls most heavily on precisely the organizations least able to absorb it. This asymmetry functions as an operational digital divide in which security and compliance capability tracks organizational wealth rather than mission need. The burden is, in this sense, regressive, a pattern documented for federal regulation generally and for cybersecurity compliance specifically, as Chapter~2 details~\parencite{Crain_Hopkins_2001, GAO_CMMC_2026}. If Digital Engineering practices prove transferable beyond defense and aerospace contexts, the capabilities that lower the labor and expertise threshold for compliance would benefit most the organizations for which that threshold is currently prohibitive. Such practices would make it more feasible for smaller organizations to demonstrate regulatory compliance and security effectiveness to funding agencies, oversight bodies, and stakeholders.

The potential for Digital Engineering to extend sophisticated security and documentation capabilities to resource-constrained organizations is a central motivating consideration for this research. Currently, enterprise architecture, authoritative traceability, and model-based documentation remain accessible primarily to large organizations with specialized expertise and dedicated budgets. Community clinics, tribal organizations, under-resourced schools, small municipalities, and social-service non-profits carry the same regulatory obligations without comparable means to meet them. The decisive question is therefore directional: whether such a capability narrows the gap between the security posture available to well-funded institutions and that available to the smallest and least-resourced organizations, or whether, like prior technology waves, it reaches the well-resourced first and widens that gap instead. This research does not directly investigate adoption within resource-constrained organizations, and it does not claim that the benefit will automatically transfer. What it contributes is the necessary first link in that chain. By establishing whether IT and Information Assurance professionals broadly perceive value in Digital Engineering capabilities, the findings create the evidentiary foundation for subsequent research into accessible, lower-cost implementations suited to the smallest and least-resourced organizations. That foundation lets the equity benefit anticipated here be deliberately pursued rather than merely hoped for.

\section{Chapter Summary}

This chapter established the context for investigating professional awareness and perceptions of Digital Engineering capabilities in Information Assurance and IT service management. It identified the challenges organizations face under traditional document-centric approaches: maintaining accurate system documentation, implementing consistent security controls, and delivering reliable IT services. Digital Engineering---with its four pillars of Model-Based Systems Engineering, digital threads, digital twin technology, and Product Lifecycle Management---offers capabilities with reported value in defense and aerospace contexts that may address analogous gaps in enterprise IT and Information Assurance practice. The institutional endorsement of Digital Engineering by the Department of War, NASA, INCOSE, and standards bodies including the Object Management Group and ISO establishes the maturity and authoritative standing of these methodologies. Enterprise architecture frameworks, particularly the Unified Architecture Framework, provide the structural foundation for implementation.

The research questions address three matters: whether professionals are aware of Digital Engineering capabilities, whether they perceive those capabilities as valuable for their work, and whether they believe the practices could improve their effectiveness in meeting compliance requirements and delivering IT services. The significance spans four dimensions. Academically, the study addresses identified literature gaps. For industry, it informs adoption strategies. Nationally, it bears on the protection of government systems. For society, it points toward better security capabilities for organizations that serve under-resourced communities.

The research targets IT and Information Assurance professionals across diverse organizational contexts, enabling assessment of awareness and perceived value across the professional community. While the defense and aerospace sectors have committed to Digital Engineering at institutional scale, this research investigates whether professionals in other sectors recognize potential value in these capabilities for their work. That question is central to the dissertation, and the chapters that follow answer it in sequence: the literature review establishes why the answer cannot be assumed, the methodology develops the instrument to measure it, and the results report what the survey finds. The findings inform whether Digital Engineering methodologies developed for defense applications might benefit organizations across all sectors, including the smallest and least-resourced.

Chapter 2 presents the literature review examining existing research across Digital Engineering, Information Assurance, and IT service management domains. The review establishes the theoretical framework for this research and documents the research gaps that this investigation begins to address. Its sustained analytical synthesis identifies what the literature contains and, beyond that, what debates remain unresolved, what tensions persist between proven capability and limited adoption, and why the present research is necessary.